% Source: https://artofproblemsolving.com/community/c3h1844637_combinations_of_dice
% Stable identity: aops-topic-1844637-q01
An outcome is described by the numbers $a_1,a_2,\dots,a_6$ of dice showing $1,2,\dots,6$, where the $a_i$ are nonnegative integers with
\[
a_1+a_2+\cdots+a_6=n .
\]
By stars and bars, write $n$ stars and $5$ bars in a row; the bars split the stars into six groups of sizes $a_1,\dots,a_6$. Every arrangement of the $n+5$ symbols gives exactly one solution and conversely, so the number of outcomes is
\[
\boxed{\dbinom{n+5}{5}=\dbinom{n+5}{n}}.
\]

The number $6^n$ counts outcomes of \emph{distinguishable} dice, where it matters which die shows which face. For identical dice, different ordered results describe the same outcome: for $n=2$, ``first die $1$, second die $2$'' and ``first die $2$, second die $1$'' are the same outcome $\{1,2\}$. So $6^2=36$ ordered results collapse to $\dbinom72=21$ outcomes ($6$ doubles and $15$ pairs of different faces). One cannot simply divide $6^n$ by $n!$ either, since outcomes with repeated faces are counted fewer than $n!$ times.
